% Part: computability
% Chapter: computability-theory
% Section: k-1

\documentclass[../../../include/open-logic-section]{subfiles}

\begin{document}

\olfileid{cmp}{thy}{k1}
\olsection{Tuladha Reduksibilitas}

Ayo kita gatekake sawijining panganggone \olref[ppr]{prop:reduce}.

\begin{prop}
\ollabel{prop:k1}
Tetepna
\[
K_1 = \Setabs{e}{\cfind{e}(0) \fdefined}.
\]
Mula $K_1$ bisa dienumerasi sacara komputabel nanging ora komputabel.
\end{prop}

\begin{proof}
Amarga $K_1 = \Setabs{e}{\lexists[s][T(e,0,s)]}$, $K_1$ bisa
dienumerasi sacara komputabel miturut \olref[eqc]{thm:exists-char}.

Kanggo nuduhake yen $K_1$ ora komputabel, ayo nuduhake yen $K_0$
bisa direduksi menyang himpunan iku.

\begin{explain}
Iki rada angel, amarga nganggo $K_1$ kita mung bisa
takon bab komputasi kang diwiwiti nganggo input tartamtu,
$0$. Upamane kowe nduwé kanca pinter kang bisa mangsuli
pitakon kaya mangkono (kanca kaya iki diarani ``orakel'').
Banjur upamane ana wong marani kowe lan takon apa
$\tuple{e,
  x}$ kalebu $K_0$, yaiku apa mesin $e$ mandheg nalika
nampa input $x$. Salah siji cara yaiku gawe mesin liyane,
$e_x$, kang kanggo input \emph{sembarang} ora nggatekake
input iku lan malah nglakokake~$e$ nganggo input $x$.
Mula cetha yen pitakon apa mesin $e$ mandheg nganggo
input $x$ ekuivalen karo pitakon apa mesin $e_x$
mandheg nganggo input $0$ (utawa input liyane apa wae).
Dadi, kowe takon marang kancamu apa mesin anyar iki,
$e_x$, mandheg nganggo input $0$; wangsulane kancamu
kanggo pitakon kang diowahi iku menehi wangsulan
kanggo pitakon asline. Iki ngasilake reduksi $K_0$
menyang~$K_1$ kang dikarepake.
\end{explain}

Kanthi nggunakake fungsi parsial komputabel universal,
tetepna $f$ minangka fungsi telung argumen kang
ditetepake dening
\[
f(x,y,z) \simeq \cfind{x}(y).
\]
Gatekna yen $f$ babar pisan ora nggatekake input katelune.
Pilihen indeks $e$ saengga $f = \cfind{e}[3]$; mula kita nduwé
\[
\cfind{e}[3](x,y,z) \simeq \cfind{x}(y).
\]
Miturut teorema $s$-$m$-$n$, ana fungsi $s(e,x,y)$ saengga,
kanggo saben $z$,
\begin{align*}
\cfind{s(e,x,y)}(z) & \simeq \cfind{e}[3](x,y,z) \\
& \simeq \cfind{x}(y).
\end{align*}

\begin{explain}
Miturut andharan informal ing ndhuwur, $s(e,x,y)$
minangka indeks kanggo mesin kang, kanggo input $z$
sembarang, ora nggatekake input iku lan ngitung
$\cfind{x}(y)$.
\end{explain}

Mligine, kita nduwé
\[
\cfind{s(e,x,y)}(0) \fdefined \quad \text{yen lan mung yen} \quad
\cfind{x}(y) \fdefined.
\]
Kanthi tembung liya, $\tuple{x, y} \in K_0$ yen lan mung yen $s(e,x,y) \in
K_1$. Mula fungsi $g$ kang ditetepake dening
\[
g(w) = s(e,(w)_0,(w)_1)
\]
minangka reduksi $K_0$ menyang~$K_1$.
\end{proof}

\end{document}
